%%
%% AIware 2026 — Benchmark & Dataset Track
%% Short Paper (4 pages incl. references)
%% Double-anonymous submission
%%
\documentclass[sigconf,anonymous]{acmart}

\usepackage{booktabs}
\usepackage{listings}
\usepackage{xcolor}
\usepackage{graphicx}
%% NOTE: do NOT add \usepackage{hyperref} — acmart loads it internally.

%% ACM metadata — fill in after acceptance
\acmConference[AIware 2026]{Proceedings of the 3rd International Conference on AI-powered Software}{April 2026}{Montreal, QC, Canada}
\acmBooktitle{AIware 2026: Proceedings of the 3rd International Conference on AI-powered Software}
\acmDOI{}
\acmISBN{}
\setcopyright{acmlicensed}

% ─────────────────────────────────────────────────────────────────────────────
\begin{document}

\title{Bridging the Data Gap: Towards Infrastructure for\\
       Custom, Dynamic Dataset Curation \& Analysis}

%% For double-anonymous: the [anonymous] class option suppresses this block
%% in the submitted PDF. Restore for camera-ready.
\author{Ademola Adeniyi}
\email{adeniyiademola18@gmail.com}
\affiliation{%
  \institution{Independent Researcher}
  \city{Philadelphia}
  \state{PA}
  \country{USA}
}

\author{Brittany Johnson}
\email{johnsonb@gmu.edu}
\affiliation{%
  \institution{George Mason University}
  \department{Department of Computer Science}
  \city{Fairfax}
  \state{VA}
  \country{USA}
}

% ─────────────────────────────────────────────────────────────────────────────
\begin{abstract}
In the age of generative artificial intelligence (AI), data is at the center of
our ability to adequately and sustainably develop advanced technologies that can
support diverse end users.
While for some populations and data types it is fairly easy to find usable,
representative datasets and sources, it becomes challenging when the focus is on
historically marginalized communities.
This makes it difficult for research like our own, which aims to study media
coverage, social discourse, and community-specific narratives to improve the
cultural alignment of generative AI technologies.
Existing news and media corpora aggregate hundreds of outlets but systematically
under- or misrepresent niche and community-centered perspectives, presenting a
gap in both existing data and dataset infrastructures.
To fill this gap, and in alignment with our ongoing efforts to support cultural
alignment of generative AI technologies, we present a modular, open-source
pipeline that continuously scrapes, normalises, and thematically annotates
articles from seven prominent Black media publications on a rolling 30-day
horizon.
The pipeline delivers approximately 1,200 articles per monthly cycle, each tagged
against a curated taxonomy of twelve social conflict topics using a hybrid
TF-IDF + seed-phrase matcher, and exports research-ready CSV artefacts alongside
a styled PDF editorial report, fully automated via GitHub Actions and deposited
to Google Drive.
We describe the pipeline architecture, dataset schema, quality metrics, and the
design decisions that make the system extensible to other under-represented
publication and data ecosystems.
\end{abstract}

\begin{CCSXML}
<ccs2012>
   <concept>
       <concept_id>10002951.10003317.10003347.10003352</concept_id>
       <concept_desc>Information systems~Data extraction and integration</concept_desc>
       <concept_significance>500</concept_significance>
   </concept>
   <concept>
       <concept_id>10002951.10003317.10003331</concept_id>
       <concept_desc>Information systems~Information retrieval</concept_desc>
       <concept_significance>300</concept_significance>
   </concept>
   <concept>
       <concept_id>10010147.10010178.10010187</concept_id>
       <concept_desc>Computing methodologies~Natural language processing</concept_desc>
       <concept_significance>300</concept_significance>
   </concept>
</ccs2012>
\end{CCSXML}

\ccsdesc[500]{Information systems~Data extraction and integration}
\ccsdesc[300]{Information systems~Information retrieval}
\ccsdesc[300]{Computing methodologies~Natural language processing}

\keywords{dataset curation, media monitoring, Black press, topic modeling,
          TF-IDF, automated pipeline, social conflict, GitHub Actions}

\maketitle

% ─────────────────────────────────────────────────────────────────────────────
\section{Introduction}

Computational research on media framing, agenda setting, and coverage equity
relies on large, timely corpora~\cite{boumans2016taking,boukes2023capturing}.
Yet the datasets most commonly used in NLP and computational journalism
research---AllTheNews, CC-News, RealNews, GDELT---are aggregated from a small
set of legacy outlets whose editorial diversity is
limited~\cite{raleigh2023political}.
Publications that serve marginalised or community-specific audiences are almost
entirely absent from these corpora.

The Black press is a historically significant and demographically distinct media
ecosystem.
Outlets such as \emph{The Root}, \emph{The Grio}, \emph{Ebony}, \emph{Essence},
\emph{NewsOne}, \emph{Capital B News}, and \emph{Blavity} collectively publish
hundreds of articles per month on politics, justice, health, culture, and
economics---often foregrounding perspectives that mainstream outlets suppress or
marginalise~\cite{gallon2021black,mellinger2020parallel}.
Despite their importance for studying how racialised communities narrate social
conflict, these outlets are largely absent from existing research corpora.

This paper presents a \emph{live, configurable, open-source pipeline} designed
to close that gap.
The system scrapes the seven publications listed above on a rolling 30-day
schedule, normalises each article to a common schema, applies a hybrid
TF-IDF + seed-phrase topic tagger across twelve social conflict topic classes,
and exports analysis-ready datasets alongside a styled research report---all
without manual intervention.
Our contributions are:

\begin{enumerate}
  \item A reproducible, schema-documented dataset of Black media articles updated
        monthly via CI/CD automation.
  \item A hybrid topic-annotation layer (TF-IDF keyword extraction + seed-phrase
        matching) extensible to arbitrary topic taxonomies.
  \item An end-to-end reference architecture that demonstrates how custom,
        under-represented media ecosystems can be continuously monitored with
        modest infrastructure costs.
\end{enumerate}

% ─────────────────────────────────────────────────────────────────────────────
\section{Background \& Motivation}

\subsection{Limitations of General-Purpose News Corpora}

Large open news corpora provide convenient baselines for NLP benchmarking but
exhibit well-documented source-diversity problems.
The GDELT project~\cite{leetaru2013gdelt} indexes millions of articles globally
but weights towards high-Alexa-rank English-language outlets.
AllTheNews~\cite{trhlik2024quantifying} and CC-News~\cite{baack2024critical}
employ similar crawl-frequency heuristics that de-prioritise low-traffic domains.
Specialised datasets---MedMCQA, LegalBench, FinancialPhraseBank---address vertical
domain gaps but have no analogue for community or ethnic press coverage.

\subsection{The Black Press as a Research Target}

The seven outlets studied here reach a combined monthly unique audience estimated
at 35--50 million visitors~\cite{cuny2022black}.
Their coverage of topics such as policing, voter access, housing, maternal
health, and reparations frequently differs in framing, sourcing, and intensity
from mainstream coverage of the same events~\cite{mjelde2023international}.
Studying these differences requires a corpus that is both timely---capturing
rapidly evolving discourse---and continuously updated, because the window of
relevance for social conflict coverage is typically measured in weeks, not years.

\subsection{Gap in Automated, Configurable Monitoring Tools}

Existing media-monitoring services (e.g., Meltwater, Cision, GDELT Analysis
Service) are proprietary, costly, and not designed to export clean research
datasets.
Open alternatives such as NewsAPI or MediaCloud provide useful APIs but do not
index the Black press at sufficient depth or recency.
No open-source, self-hosted tool exists that (a) targets this specific
publication set, (b) applies a sociologically grounded topic taxonomy, and
(c) runs fully automated on free CI/CD infrastructure.

% ─────────────────────────────────────────────────────────────────────────────
\section{System Architecture \& Approach}

The pipeline consists of four decoupled stages executed sequentially by a single
orchestrator (\texttt{run.py}): \textbf{Scrape} $\to$ \textbf{Analyse} $\to$
\textbf{Report} $\to$ \textbf{Deliver}.
Each stage reads from and writes to the \texttt{outputs/} directory, enabling
any stage to be re-run independently.

\subsection{Stage 1: Scraping}

Each of the seven publications requires a distinct scraping strategy due to
differences in technology stack (Table~\ref{tab:sources}).

\begin{table}[h]
\caption{Per-publication scrape methods}
\label{tab:sources}
\small
\begin{tabular}{lll}
\toprule
\textbf{Publication} & \textbf{Method} & \textbf{Fallback} \\
\midrule
thegrio.com      & WP REST API       & RSS  \\
theroot.com      & WP REST API       & RSS  \\
newsone.com      & WP REST API       & RSS  \\
ebony.com        & WP REST API       & RSS  \\
capitalbnews.org & Sitemap + HTML    & RSS  \\
essence.com      & Paginated HTML    & RSS  \\
blavity.com      & Playwright + HTML & RSS  \\
\bottomrule
\end{tabular}
\end{table}

\noindent\textbf{WordPress REST API.}
Four outlets run WordPress and expose the \texttt{/wp-json/wp/v2/posts}
endpoint.
We issue paginated requests with \texttt{after} and \texttt{before} date
filters covering the rolling 30-day window and embed term and author metadata
in a single round trip via \texttt{\_embed=wp:term,wp:author}.

\noindent\textbf{Sitemap + HTML (Capital B News).}
Capital B's sitemap index exposes post, article, and news sub-sitemaps with
\texttt{<lastmod>} timestamps.
We parse these first (fast pre-filter), then discover additional URLs via RSS
feed, category pages (up to 3 pages), and WordPress date archives.
Candidates are capped at 200 to keep runtime under 5 minutes; sitemap entries
with \texttt{lastmod} are prioritised over undated discoveries.

\noindent\textbf{Playwright headless Chromium (Blavity).}
Blavity is a JavaScript-rendered single-page application; its sitemap URL
returns HTML rather than XML, and its RSS feed is capped at 10 items.
We use Playwright to navigate eight category pages (up to 10 pagination steps
each), intercept article links, and extract metadata including JSON-LD
structured data for publication dates.
A requests-based HTML fallback is available if Playwright is unavailable.

All scrapers share a common normalisation step and a polite crawl delay of
250--500\,ms between requests.
A unified RSS fallback triggers when any primary method returns fewer than
five articles.

\subsection{Stage 2: Analysis}

\noindent\textbf{TF-IDF Keyword Extraction.}
We concatenate the title and description of every article and vectorise the
corpus using scikit-learn's \texttt{TfidfVectorizer} with $n$-gram range
$(1,3)$, sublinear TF scaling, and NLTK English stopwords augmented with a
domain-specific extension list.
The top-$k$ terms by aggregate TF-IDF weight are reported as the period's
``signal keywords''---phrases that are both frequent and distinctive relative to
a generic background corpus.

\noindent\textbf{Seed-Phrase Topic Matching.}
We maintain a curated taxonomy of twelve social conflict topics (see
Table~\ref{tab:topics}).
Each topic is defined by a name, a one-sentence description, and a list of seed
phrases.
Matching is performed with whole-word regular expressions that tolerate common
morphological variants (e.g., \emph{redlin} matches \emph{redlining},
\emph{redlined}, \emph{redlines}).
An article is assigned to a topic if at least one seed phrase matches in the
combined title + description field.
Multiple topic assignments are permitted; the empty string is used when no
match is found.

\begin{table}[h]
\caption{Social conflict topic taxonomy (12 topics)}
\label{tab:topics}
\small
\begin{tabular}{ll}
\toprule
\textbf{Topic} & \textbf{Example seed phrases} \\
\midrule
Policing \& Public Safety     & police brutality, use of force, qualified immunity \\
Voter Suppression             & voter id, gerrymandering, polling place closures \\
Book Bans \& Anti-DEI         & book ban, critical race theory, dei ban \\
Housing \& Displacement       & eviction, gentrification, affordable housing \\
Maternal Health               & maternal mortality, birth equity, obstetric racism \\
Redlining \& Fair Housing     & redlining, discriminatory lending, fair housing \\
Anti-Black Surveillance       & facial recognition, predictive policing, surveillance \\
Reparations                   & reparations, slavery redress, atonement fund \\
School Funding                & school funding, education funding, title i \\
Criminal Justice Reform       & mass incarceration, mandatory minimum, bail reform \\
Environmental Justice         & environmental racism, cancer alley, sacrifice zone \\
Economic Equity \& Wealth Gap & racial wealth gap, predatory lending, banking desert \\
\bottomrule
\end{tabular}
\end{table}

\subsection{Stage 3: Report Generation}

A Jinja2~\cite{pallets2025jinja} HTML template is populated with the analysis
outputs and rendered to PDF via WeasyPrint~\cite{courtbouillon2021weasyprint}.
The report includes: (i) a per-source article count and data quality table,
(ii) a topic frequency bar chart, (iii) a source-by-topic heatmap,
(iv) per-topic TF-IDF keyword word clouds, and (v) an annotated article index
with clickable URLs grouped by topic.
Chart rendering uses matplotlib and seaborn with a white background to satisfy
academic publication legibility norms.

\subsection{Stage 4: Delivery \& Automation}

Output files (combined CSV, seven per-source CSVs, HTML report, PDF report) are
uploaded to a designated Google Drive folder via the Drive v3 REST API with
OAuth 2.0 refresh-token authentication.
Existing files are updated in-place rather than duplicated to keep the folder
stable for downstream consumers.

The pipeline is fully automated via GitHub Actions (ubuntu-22.04), running on
the first of every month via a scheduled \texttt{cron} trigger.
It also supports a \texttt{workflow\_dispatch} event, which allows any
collaborator with repository access to trigger a full pipeline run manually
at any time---useful for reprocessing after a scraper fix, testing a new topic
addition, or producing an ad-hoc report outside the regular monthly cadence.
End-to-end wall-clock time is approximately 5--8 minutes; Playwright
installation and Chromium download account for roughly 90 seconds of that.

% ─────────────────────────────────────────────────────────────────────────────
\section{Dataset Characteristics}

\subsection{Coverage}

In testing over the February--March 2026 window, the pipeline collected
1,218 articles across the seven sources (Table~\ref{tab:coverage}).
Article counts per source vary significantly, reflecting both publication
cadence and the diversity of scraping methods required.

\begin{table}[h]
\caption{Approximate monthly article yield per source (Feb--Mar 2026)}
\label{tab:coverage}
\small
\begin{tabular}{lrrrr}
\toprule
\textbf{Source} & \textbf{Articles} & \textbf{Date\%} & \textbf{Snippet\%} & \textbf{Cat.\%} \\
\midrule
thegrio.com      & 389 & 100 & 97 & 98 \\
theroot.com      & 312 & 100 & 95 & 95 \\
newsone.com      & 278 & 100 & 93 & 91 \\
ebony.com        &  84 & 100 & 88 & 72 \\
essence.com      &  71 &  96 & 82 & 68 \\
capitalbnews.org &  51 & 100 & 91 & 85 \\
blavity.com      &  33 &  88 & 79 & 61 \\
\midrule
\textbf{Total}   & \textbf{1,218} & \textbf{99} & \textbf{93} & \textbf{90} \\
\bottomrule
\end{tabular}
\end{table}

\subsection{Schema}

Each article record contains the fields in Table~\ref{tab:schema}.
The \texttt{topics} field appears only in \texttt{articles\_tagged.csv}.

\begin{table}[h]
\caption{Dataset schema}
\label{tab:schema}
\small
\begin{tabular}{lll}
\toprule
\textbf{Field} & \textbf{Type} & \textbf{Description} \\
\midrule
title               & string & Article headline \\
description         & string & Excerpt / meta description \\
source              & string & Publication domain \\
date\_of\_publication & date   & YYYY-MM-DD \\
category            & string & Publication-assigned category \\
author              & string & Byline \\
link                & string & Canonical URL \\
topics              & string & Matched topic(s) or empty \\
\bottomrule
\end{tabular}
\end{table}

\subsection{Topic Distribution}

Across the test window, \textbf{Policing \& Public Safety} (18\%),
\textbf{Criminal Justice Reform} (14\%), and \textbf{Economic Equity \&
Wealth Gap} (12\%) were the most frequently matched topics.
\textbf{Reparations} (3\%) and \textbf{Anti-Black Surveillance} (4\%) had the
lowest match rates, consistent with their more targeted seed-phrase sets.
Approximately 31\% of articles received no topic match, indicating
non-conflict-focused content (culture, sports, entertainment) or articles whose
match tokens appear in the body text but not in the title/description field.
These distributions are consistent with prior characterisations of Black press
editorial priorities~\cite{mjelde2023international,gallon2021black}.

% ─────────────────────────────────────────────────────────────────────────────
\section{Limitations \& Future Work}

\subsection{Scraping Reliability}

\noindent\textbf{JavaScript-rendered sites require brittle workarounds.}
Three of the seven outlets---Blavity, Essence, and Capital B News---do not expose
clean server-rendered HTML or functional WordPress REST APIs, requiring
specialised scraping strategies that are inherently more fragile than API-based
approaches.
Blavity, a fully JavaScript-rendered single-page application, presented the
most acute challenges.
Its sitemap endpoint returns an HTML shell rather than XML, making standard
sitemap discovery impossible.
Publication dates are injected into the DOM by client-side JavaScript and are
absent from raw server responses; the scraper recovers dates only by reading
\texttt{application/ld+json} structured-data blocks embedded in the initial
HTML payload.
Internal links on category listing pages use root-relative paths
(e.g., \texttt{/article-slug}) rather than absolute URLs, requiring explicit
base-URL prepension before domain filtering.
Category URLs follow the pattern \texttt{/categories/X} rather than the
conventional \texttt{/category/X}; requests to the latter return HTTP 404.
Each of these issues independently caused the Blavity scraper to return zero
valid articles during development, and diagnosing them required iterative
live testing against the production site.

\noindent\textbf{Anti-crawling measures suppress yield.}
Even with a Playwright headless Chromium browser, Blavity's CDN occasionally
serves challenge pages to automated clients, reducing the scraper's effective
yield to as few as one article per run before the RSS fallback fires.
Capital B News, which publishes approximately 28 articles per month, exposes
over 1,000 candidate URLs through its sitemap and category pages; naively
scraping all of them increases per-run wall-clock time from under 5 minutes to
over 10 minutes.
The current implementation caps scraping candidates at 200 and limits category
pagination to 3 pages, which recovers the full recent window at the cost of
potentially missing older articles near the 30-day boundary.

\noindent\textbf{Scraper fragility under site redesigns.}
HTML-based scrapers are coupled to the DOM structure of each target site.
A site redesign---changed CSS selectors, restructured navigation, or a migration
away from WordPress---will silently reduce article yield without raising an
exception.
The WordPress REST API scrapers (The Grio, The Root, NewsOne, Ebony) are
comparatively robust and require no structural DOM assumptions.
We plan to add per-source yield regression tests on a scheduled basis that
alert when any source returns significantly fewer articles than its historical
baseline, distinguishing scraper breakage from genuine low-publication periods.

\subsection{Coverage Asymmetry}

Article yield varies substantially across sources: high-cadence outlets such as
The Grio and The Root contribute approximately 300--400 articles per month each,
while Blavity contributes approximately 30--35 and Capital B News approximately
50~\cite{cuny2022black}.
This asymmetry means that aggregate TF-IDF scores and topic frequency counts
are dominated by the vocabulary and coverage priorities of higher-volume outlets.
Researchers using this dataset for cross-outlet comparative analysis should
weight or stratify by source before drawing conclusions about community-level
discourse.

\subsection{Topic Annotation Quality}

\noindent\textbf{Surface-text matching only.}
The seed-phrase matcher operates exclusively on the article title and
meta-description fields.
Approximately 31\% of articles in our test window received no topic match,
reflecting both genuinely off-topic content (entertainment, sports, culture) and
cases where the relevant topic terms appear only in the article body and not in
the title or description.
Full body-text matching would substantially improve recall for the latter group.
A planned extension will optionally store and match against extracted body text.

\noindent\textbf{Label boundary overlap.}
Seed-phrase matching produces noisy annotations at the intersection of related
topics.
A story about school segregation in a historically redlined neighbourhood may
legitimately match both \emph{School Funding} and \emph{Redlining \& Fair
Housing}; a story about disproportionate incarceration of Black men may match
both \emph{Criminal Justice Reform} and \emph{Economic Equity \& Wealth Gap}.
The pipeline permits multi-label assignment, but the seed-phrase sets were
defined independently and do not encode inter-topic precedence or exclusion
rules.
Future work will evaluate zero-shot classification using large language models
with the topic descriptions as natural-language hypotheses
(e.g., via a Natural Language Inference head), producing confidence-calibrated
multi-label annotations that better reflect editorial nuance.

\noindent\textbf{Towards a fully agentic pipeline.}
Beyond improving individual annotation steps, we plan to develop an end-to-end
agentic process in which LLMs drive both data collection and analysis.
In this envisioned architecture, an LLM agent would (i) discover new
publication sources relevant to a user-specified community or topic domain,
(ii) generate and iteratively refine scraper code for those sources,
(iii) synthesise topic ontologies grounded in the collected corpus rather than
pre-specified seed phrases, and (iv) produce natural-language analytical
summaries alongside the structured dataset outputs.
This would transform the pipeline from a fixed-topology ETL system into a
self-configuring research assistant capable of targeting arbitrary
underrepresented media ecosystems with minimal human intervention.

\noindent\textbf{No sentiment or framing analysis.}
The current pipeline identifies \emph{what} topics are covered but not
\emph{how} they are framed.
Two articles both matched to \emph{Policing \& Public Safety} may take
diametrically opposed stances---one celebrating a reform, one decrying police
defunding.
Extending the annotation layer with sentiment scores or frame classifiers
(e.g., using a media frames corpus as a training signal) would substantially
increase the dataset's utility for framing studies.

\subsection{Temporal and Structural Constraints}

\noindent\textbf{Rolling 30-day window.}
Each pipeline run retains only articles published within the preceding 30 days.
Longitudinal analysis spanning multiple months requires aggregating the
versioned outputs across runs---now supported by the dated-subfolder output
structure, but not automatically merged.
A future pipeline stage will maintain a deduplicated cumulative corpus alongside
the per-run snapshots.

\noindent\textbf{No article body storage.}
The schema stores only titles and meta-descriptions, not full article text,
to minimise storage overhead and avoid copyright concerns around bulk
reproduction of publisher content.
This limits the dataset's utility for tasks that require body-level features
(e.g., readability scoring, named entity recognition, full-text topic modeling).
We are exploring a lightweight opt-in body extraction mode gated behind
explicit researcher acknowledgement of fair-use terms.

\subsection{Extensibility to Other Ecosystems}

The pipeline's modular design---each source has an isolated scraper module, a
shared schema normaliser, and a declarative entry in the publication registry---
was deliberately chosen to support extension to other under-represented
publication sets (e.g., Indigenous press, Spanish-language outlets, LGBTQ+
media, diaspora publications).
The topic taxonomy is fully configurable through a single file
(\texttt{analysis/topics.py}) with no changes required elsewhere.
However, onboarding a new outlet requires manual investigation of its technical
stack (WordPress vs.\ SPA vs.\ static site), its URL structure, its date
encoding conventions, and whether it is accessible to automated clients---a
process that took between 30 minutes (WordPress outlets) and several days
(Blavity) during the development of the current system.

% ─────────────────────────────────────────────────────────────────────────────
\section{Conclusion}

We have described an open-source pipeline that addresses a concrete gap in
computational journalism and social-computing research: the absence of
continuously updated, community-press datasets annotated against sociologically
grounded topic taxonomies.
The system delivers approximately 1,200 research-ready articles per monthly
cycle from seven Black media outlets, annotated across twelve social conflict
topics, with no recurring manual intervention.
We believe the architecture---modular scrapers, declarative publication registry,
extensible topic taxonomy, CI/CD scheduling---serves as a replicable template
for researchers who need live, domain-specific datasets for communities whose
media presence is currently invisible to mainstream corpora.
Source code, dataset samples, and documentation are available at
\texttt{[repository URL redacted for review]}.

% ─────────────────────────────────────────────────────────────────────────────
\begin{acks}
Omitted for anonymous review.
\end{acks}

% ─────────────────────────────────────────────────────────────────────────────
\bibliographystyle{ACM-Reference-Format}
\begin{thebibliography}{12}

\bibitem{boumans2016taking}
Jelle W. Boumans and Damian Trilling. 2016.
Taking Stock of the Toolkit: An Overview of Relevant Automated Content Analysis
Approaches and Techniques for Digital Journalism Scholars.
\textit{Digital Journalism} 4, 1 (2016), 8--23.
DOI: \url{https://doi.org/10.1080/21670811.2015.1096598}

\bibitem{boukes2023capturing}
Mark Boukes and Damian Trilling. 2023.
Capturing a News Frame: Comparing Machine-Learning Approaches to Frame Analysis
with Different Degrees of Supervision.
\textit{Communication Methods and Measures} 17, 4 (2023), 347--371.
DOI: \url{https://doi.org/10.1080/19312458.2023.2230560}

\bibitem{raleigh2023political}
Clionadh Raleigh, Roudabeh Kishi, and Andrew Linke. 2023.
Political Instability Patterns Are Obscured by Conflict Dataset Scope Conditions,
Sources, and Coding Choices.
\textit{Humanities and Social Sciences Communications} 10 (2023), 74.
DOI: \url{https://doi.org/10.1057/s41599-023-01559-4}

\bibitem{gallon2021black}
Kim Gallon. 2021.
The Black Press.
In \textit{Oxford Research Encyclopedia of American History}.
Oxford University Press.
DOI: \url{https://doi.org/10.1093/acrefore/9780199329175.013.851}

\bibitem{mellinger2020parallel}
Gwyneth Mellinger. 2020.
A Parallel Reality: The Challenges of Black Press History.
\textit{Journalism \& Mass Communication Quarterly} 98, 1 (2020), 8--30.
DOI: \url{https://doi.org/10.1177/1522637920965960}

\bibitem{leetaru2013gdelt}
Kalev Leetaru and Philip A. Schrodt. 2013.
GDELT: Global Data on Events, Location, and Tone, 1979--2012.
In \textit{Proceedings of the International Studies Association Annual Convention}.
\url{http://data.gdeltproject.org/documentation/ISA.2013.GDELT.pdf}

\bibitem{trhlik2024quantifying}
Filip Trhlik and Pontus Stenetorp. 2024.
Quantifying Generative Media Bias with a Corpus of Real-World and Generated
News Articles.
In \textit{Findings of ACL: EMNLP 2024}, 4420--4445.
DOI: \url{https://doi.org/10.18653/v1/2024.findings-emnlp.255}

\bibitem{baack2024critical}
Stefan Baack. 2024.
A Critical Analysis of the Largest Source for Generative AI Training Data:
Common Crawl.
In \textit{Proceedings of the 2024 ACM FAccT Conference}, 2199--2208.
DOI: \url{https://doi.org/10.1145/3630106.3659033}

\bibitem{cuny2022black}
CUNY Craig Newmark Graduate School of Journalism. 2022.
Black Media Power: Connecting Us to the Historical Fight for Justice.
Center for Community Media, New York, NY.
Retrieved March 28, 2026 from \url{https://blackmediareport.journalism.cuny.edu/connecting-us-to-the-historical-fight-for-justice/}

\bibitem{mjelde2023international}
Hilmar Mjelde. 2023.
International News Media Coverage of Black Lives Matter: Evidence From Norway.
\textit{International Journal of Communication} 17 (2023), 3679--3699.
\url{https://ijoc.org/index.php/ijoc/article/view/19490}

\bibitem{courtbouillon2021weasyprint}
CourtBouillon. 2021.
WeasyPrint: The Awesome Document Factory (Version 53.3).
\url{https://github.com/Kozea/WeasyPrint}

\bibitem{pallets2025jinja}
Pallets. 2025.
Jinja: A Fast, Expressive, Extensible Templating Engine (Version 3.1.6).
Retrieved March 29, 2026 from \url{https://jinja.palletsprojects.com/}

\end{thebibliography}

\end{document}
